\section{Stars: the cube condition for \texorpdfstring{$\SU(2)$}{SU(2)} past \texorpdfstring{$\beta_W=\frac13$}{beta\_W=1/3}}\label{sec:stars}

Proposition~\ref{prop:methodlimit} shows that no argument built on single-link
updates and gauge-fixed links passes $\beta_W=\frac13$ for $\SU(2)$. This
section replaces single links by \emph{stars}, the eight links at a vertex of
one parity, and passes that value: the cube condition holds uniformly in the
cube size for $|\beta_W|\le\frac7{20}$
(Theorem~\ref{thm:starsuniform}), and a block condition giving a lattice gap
holds for $|\beta_W|\le\frac9{25}$ (Theorem~\ref{thm:stars}). The gain is
visible at first order: the perturbation of one plaquette at a vertex enters
the law of the star only through the relative variable of two of its links,
so a star responds with the coefficient $c$ where two single links respond
with $2c$, and no plaquette is internal to a star.\footnote{A Dobrushin
certificate bounds the response of a block to a change of its boundary by a
matrix of local influences, and contraction means that this matrix has
spectral radius below one. Grouping links into blocks changes the matrix:
interactions inside a block are treated exactly instead of being bounded.}

\emph{Setting.} Throughout the section $G=\SU(2)$, identified with the unit
quaternions $S^3\subset\R^4$, with the Euclidean inner product
$\ip\cdot\cdot$ of $\R^4$ and the chordal metric $\delta(U,V)=|U-V|$, which is
bi-invariant, invariant under inversion, and of diameter $2$. For unit
quaternions $\ReTr U=2\ip U1$, and for unit $p,q$ and $x,y\in\R^4$,
$\ip{px}y=\ip x{\bar py}$, $\ip{xq}y=\ip x{y\bar q}$ and $|pxq|=|x|$. The
Wilson measure is $\mu(dU)=Z^{-1}\prod_pe^{2\beta\ReTr U_p}\prod_edU_e$ on
$\Lambda_L$, $L$ even, with $\beta_W=4\beta$, and $\mathcal A$, $V_x$,
$\mathcal P_x$ and $\BP(n,\varpi)$ are as in Sec.~\ref{sec:cube}. For a set
$V$ of links, $\mathcal P_V$ is the conditional expectation given the links
outside $V$. For $b\in\R^4$, $\nu_b(dx)\propto e^{\ip xb}dx$ is the von
Mises--Fisher law on $S^3$, and the sharp quarter (Theorem~\ref{q-thm:quarter}
in Appendix~\ref{app:sharp}) states
\begin{equation}\label{eq:Q}
W_1(\nu_b,\nu_{b'})\le\theta\,|b-b'|\quad\text{for all }b,b'\in\R^4,
\qquad\theta=\tfrac14,
\end{equation}
with $W_1$ the Kantorovich distance for $\delta$. This is the only property of
$\nu_b$ used in the section. We put
\[
c=4\theta|\beta|=|\beta|=\tfrac14|\beta_W| .
\]
Every argument below uses only~\eqref{eq:Q}, and with a constant $\theta$ in
place of $\frac14$ the ranges become $4\theta|\beta|\le\frac7{80}$ and
$4\theta|\beta|\le\frac9{100}$.

\subsection{Even stars and two-paths}

Call a vertex \emph{even} if the sum of its coordinates is even. For an even
vertex $v$ the \emph{star} $B_v$ is the set of the eight links with an end
point at $v$. Write the eight directions as $\pm1,\dots,\pm4$, and for the link
of $B_v$ in direction $\sigma$ let $X_\sigma$ be its parallel transporter from
$v$ outwards: $X_{+\mu}=U_{(v,\mu)}$ and $X_{-\mu}=U_{(v-e_\mu,\mu)}^{-1}$. Two
directions are \emph{adjacent} if they are distinct and not opposite; each
direction is adjacent to six others, and there are $24$ adjacent pairs.

\begin{lemma}\label{lem:stars}
Let $L\ge4$ be even.
\begin{enumerate}[label=\textup{(\roman*)},leftmargin=*]
\item Every link has exactly one even end point, so the stars partition $E$.
\item Every plaquette has exactly two even corners $v,v'$, which are opposite
corners, and two of its links lie in $B_v$, two in $B_{v'}$. The plaquettes
with corner $v$ correspond bijectively to the adjacent pairs $\{\sigma,\tau\}$:
the plaquette $p_{\sigma\tau}$ has corners $v$, $v+\sigma$, $v+\sigma+\tau$,
$v+\tau$.
\item For such a plaquette, $\ReTr U_p=2\ip{h_p}{X_\sigma^{-1}X_\tau}$, where
the \emph{two-path} $h_p$ is the transporter from $v+\sigma$ to $v+\tau$
through the opposite even corner $v'=v+\sigma+\tau$, the product of the two
links of $p$ in $B_{v'}$, suitably inverted. $h_p$ is unchanged by a gauge
transformation at $v'$, and changing one of its two links by $\eta$ in
$\delta$ changes $h_p$ by exactly $\eta$.
\item The conditional law of $(X_\sigma)_\sigma$ given the links outside $B_v$ is
\[
\gamma_{v,h}(dX)=\frac1{Z_v}\prod_{\{\sigma,\tau\}}e^{4\beta\ip{h_{\sigma\tau}}
{X_\sigma^{-1}X_\tau}}\prod_\sigma dX_\sigma ,
\]
with one ordering fixed in each adjacent pair. It depends on the
configuration only through the $24$ two-paths $h=(h_{\sigma\tau})$, and it is
invariant under $X_\sigma\mapsto gX_\sigma$, the gauge transformation at $v$.
\item A link $i$ with even end point $v_i$ lies in exactly six plaquettes;
their even corners other than $v_i$ are six distinct even vertices, and for
each of them, $v$, the link $i$ is a factor of exactly one two-path $h_p$ with
$p$ at $v$.
\end{enumerate}
\end{lemma}

\begin{proof}
(i) The end points of a link differ by a unit vector, so their parities
differ. (ii) The corners of a plaquette are $y$, $y+e_\mu$, $y+e_\mu+e_\nu$,
$y+e_\nu$, of parities $\epsilon,\epsilon+1,\epsilon,\epsilon+1$, and each
link joins consecutive corners. A plaquette with even corner $v$ is spanned by
two adjacent directions at $v$, and conversely; since $L\ge4$ these $24$
squares are distinct. (iii) Going around $p$ from $v$, $U_p$ is conjugate to
$X_\sigma h_pX_\tau^{-1}$ or to its inverse, and $\ReTr$ is a class function,
invariant under inversion on $\SU(2)$; so $\ReTr U_p=\ReTr(X_\sigma h_p
X_\tau^{-1})=2\ip{X_\sigma h_p}{X_\tau}=2\ip{h_p}{X_\sigma^{-1}X_\tau}$. A gauge
transformation $g$ at $v'$ multiplies the first factor of $h_p$ on the right
by $g^{-1}$ and the second on the left by $g$, and multiplication by unit
quaternions is an isometry. (iv) The plaquettes containing a link of $B_v$
are those with corner $v$, by (ii), and their weights are given by (iii); the
gauge transformation at $v$ changes no two-path. (v) If $i$ joins $v_i$ to
$w=v_i+\sigma$, its plaquettes are spanned by $\sigma$ and the six directions
$\tau\perp\sigma$, with opposite even corners $v_i+\sigma+\tau$, which are
distinct since $L\ge4$; in the plaquette with corner $v$, the link $i$ is one
of the two links at $v_i$, hence a factor of the two-path.
\end{proof}

For the star at $v$ define the quotient pseudometric
\[
\bar d_v(X,X')=\min_{g\in G}\sum_\sigma\delta(gX_\sigma,X'_\sigma).
\]

\subsection{The seven-spin system}

Fix an adjacent pair $\{a,b\}$, let $S_a$ be the set of the seven directions
other than $a$, $A_7$ its adjacency matrix and $e_b$ the indicator of $b$.
In $S_a$, the direction $b$ is adjacent to $\bar a=-a$ and to the four
directions $C$ of the other two axes; $\bar a$ is adjacent to $b$,
$\bar b=-b$ and $C$; $\bar b$ to $\bar a$ and $C$; and each element of $C$ to
$b$, $\bar a$, $\bar b$ and the two elements of $C$ on the other axis. The
largest degree is $6$. For $0\le c<\frac16$ put
\[
s(c)=\mathbf 1^{\mathsf T}(I-cA_7)^{-1}e_b=\sum_{k\ge0}c^k\,\mathbf 1^{\mathsf T}
A_7^ke_b ,
\]
a nondecreasing function of $c$ with $s(c)\ge1$. The symmetries fixing $a$ and
$b$ permute $C$ transitively, so $x=(I-cA_7)^{-1}e_b$ is constant, equal to
$z$, on $C$, and $(x_b,x_{\bar a},x_{\bar b},z)$ is the unique solution of
\begin{equation}\label{eq:sevensys}
\begin{aligned}
x_b&=1+c(x_{\bar a}+4z),&\quad x_{\bar a}&=c(x_b+x_{\bar b}+4z),\\
x_{\bar b}&=c(x_{\bar a}+4z),&\quad z&=c(x_b+x_{\bar a}+x_{\bar b}+2z),
\end{aligned}
\end{equation}
with $s(c)=x_b+x_{\bar a}+x_{\bar b}+4z$.

\begin{lemma}\label{lem:svalues}
\begin{enumerate}[label=\textup{(\roman*)},leftmargin=*]
\item $s(\frac1{12})=\frac{26}{15}$, with $(x_b,x_{\bar a},x_{\bar b},z)=
\frac1{210}(221,28,11,26)$.
\item $s(\frac7{80})=\frac{3480}{1933}$, with $(x_b,x_{\bar a},x_{\bar b},z)=
\frac1{181702}(192531,26320,10829,24360)$.
\item $s(\frac9{100})=\frac{5450}{2957}$, with $(x_b,x_{\bar a},x_{\bar b},z)=
\frac1{348926}(371363,53100,22437,49050)$.
\end{enumerate}
Hence $6c\,s(c)$ equals $\frac{13}{15}$ at $c=\frac1{12}$, $\frac{1827}{1933}$
at $c=\frac7{80}$ and $\frac{2943}{2957}$ at $c=\frac9{100}$, and
$6c\,s(c)\le\frac{2943}{2957}$ for $0\le c\le\frac9{100}$.
\end{lemma}

\begin{proof}
Since $c\lambda_{\max}(A_7)\le6c<1$, the matrix $I-cA_7$ is invertible and it
suffices to check~\eqref{eq:sevensys}. (i) Multiplied by $12\cdot210$ the four
equations read $12\cdot221-(28+104)=2520$, $12\cdot28-(221+11+104)=0$,
$12\cdot11-132=0$ and $12\cdot26-(221+28+11+52)=0$; the sum is
$364/210=26/15$. (ii) Multiplied by $80\cdot181702=14536160$:
\begin{align*}
80\cdot192531-7(26320+4\cdot24360)&=15402480-866320=14536160,\\
80\cdot26320-7(192531+10829+4\cdot24360)&=2105600-7\cdot300800=0,\\
80\cdot10829-7(26320+4\cdot24360)&=866320-7\cdot123760=0,\\
80\cdot24360-7(192531+26320+10829+2\cdot24360)&=1948800-7\cdot278400=0.
\end{align*}
The sum is $327120/181702=3480/1933$, as $327120=94\cdot3480$ and
$181702=94\cdot1933$. (iii) Multiplied by $100\cdot348926$:
$100\cdot371363-9(53100+196200)=34892600$,
$100\cdot53100-9(371363+22437+196200)=5310000-9\cdot590000=0$,
$100\cdot22437-9(53100+196200)=0$ and
$100\cdot49050-9(371363+53100+22437+98100)=4905000-9\cdot545000=0$. The sum
is $643100/348926=5450/2957$, as $348926=118\cdot2957$ and
$643100=118\cdot5450$. The last claim follows from the monotonicity of $s$.
\end{proof}

\begin{lemma}[one star]\label{lem:onestar}
Assume~\eqref{eq:Q} and $c<\frac16$. Let $h,h'$ be two families of two-paths
at $v$ that differ only at one adjacent pair $\{a,b\}$. Then
\[
W_1^{\bar d_v}(\gamma_{v,h},\gamma_{v,h'})\le c\,s(c)\,|h_{ab}-h'_{ab}| .
\]
\end{lemma}

\begin{proof}
Under $\gamma_{v,h}$ put $Z_\sigma=X_a^{-1}X_\sigma$. The map $X\mapsto(X_a,Z)$
preserves Haar measure, and by Lemma~\ref{lem:stars}(iv) the density depends
only on $Z$, with $Z_a=1$; so $X_a$ is Haar distributed and independent of
$Z_S=(Z_\sigma)_{\sigma\in S_a}$, whose law $\mu^{(a)}_h$ has density
proportional to $\prod e^{4\beta\ip{h_{\sigma\tau}}{Z_\sigma^{-1}Z_\tau}}$. Given
a coupling of $\mu^{(a)}_h$ and $\mu^{(a)}_{h'}$, take $X_a=X'_a$ Haar and
independent, $X_\sigma=X_aZ_\sigma$ and $X'_\sigma=X_aZ'_\sigma$; this couples
$\gamma_{v,h}$ and $\gamma_{v,h'}$, and with $g=1$ in the definition,
$\bar d_v(X,X')\le\sum_{\sigma\in S_a}\delta(Z_\sigma,Z'_\sigma)$. It suffices
to bound the Kantorovich distance of $\mu^{(a)}_h$ and $\mu^{(a)}_{h'}$ for the
cost $\sum_{\sigma\in S_a}\delta$.

For $\sigma\in S_a$ the conditional law of $Z_\sigma$ given the others is
$\nu_{B_\sigma}$ with $B_\sigma=4\beta\sum_{\tau\text{ adj.\ }\sigma}w_{\sigma\tau}$,
where $w_{\sigma\tau}=Z_\tau\bar h_{\sigma\tau}$ or $Z_\tau h_{\tau\sigma}$
according to the ordering of the pair, since $\ip{h_{\sigma\tau}}{Z_\sigma^{-1}
Z_\tau}=\ip{Z_\sigma}{Z_\tau\bar h_{\sigma\tau}}$; changing $Z_\tau$ by $\eta$
changes $w_{\sigma\tau}$ by $\eta$, and changing $h_{ab}$ changes only $B_b$, by
at most $4|\beta||h_{ab}-h'_{ab}|$. Let $\pi$ minimise
$\int\sum_\sigma\delta(Z_\sigma,Z'_\sigma)\,d\pi$ over couplings of
$\mu^{(a)}_h$ and $\mu^{(a)}_{h'}$ (one exists by compactness) and put
$\varrho_\sigma=\int\delta(Z_\sigma,Z'_\sigma)\,d\pi$. Replacing
$(Z_\sigma,Z'_\sigma)$ by a measurably chosen optimal coupling of the two
conditional laws~\cite[Cor.~5.22]{Vil09} gives another coupling that changes
only the $\sigma$-term of the cost, so by minimality and~\eqref{eq:Q}
\[
\varrho_\sigma\le\theta\int|B_\sigma-B'_\sigma|\,d\pi\le c\sum_{\tau\in S_a,\
\tau\text{ adj.\ }\sigma}\varrho_\tau+c\,|h_{ab}-h'_{ab}|\,[\sigma=b],
\]
that is, $\varrho\le cA_7\varrho+c|h_{ab}-h'_{ab}|e_b$ componentwise. Iterating,
$\varrho\le\sum_{k<K}(cA_7)^kc|h_{ab}-h'_{ab}|e_b+(cA_7)^K\varrho$, and
$(cA_7)^K\to0$ since the spectral radius of $cA_7$ is at most $6c<1$.
Summing the components gives $\sum_\sigma\varrho_\sigma\le c\,s(c)|h_{ab}-
h'_{ab}|$.
\end{proof}

\subsection{The star condition and a lattice gap up to \texorpdfstring{$\beta_W=\frac9{25}$}{beta\_W=9/25}}

Let $\mathcal P_v=\mathcal P_{B_v}$ and let $N_s=L^4/2$ be the number of even
vertices. For a function $h$ on $G^E$ let $\Lip(h)$ be the least constant
with $|h(U)-h(U')|\le\Lip(h)\delta(U_i,U'_i)$ whenever $U,U'$ differ only at
one link $i$; changing links one at a time, $|h(U)-h(U')|\le\Lip(h)\sum_e
\delta(U_e,U'_e)$.

\begin{proposition}[star condition]\label{prop:starcond}
Assume~\eqref{eq:Q} and put $\gamma=c\,s(c)$, with $c<\frac16$ and $6\gamma<1$.
Then for every real $f\in\mathcal A$
\[
\sum_{v\text{ even}}\mu\bigl((f-\mathcal P_vf)^2\bigr)\ge(1-6\gamma)\Var_\mu(f).
\]
\end{proposition}

\begin{proof}
Let $h\in\mathcal A$ be Lipschitz, gauge invariant at every configuration, and
let $U,U'$ differ only at $i$. If $i\in B_v$, then $\mathcal P_vh(U)=\mathcal
P_vh(U')$. Otherwise let $\pi$ be a coupling of $\gamma_v(\cdot\mid U)$ and
$\gamma_v(\cdot\mid U')$ minimising the integral of $\bar d_v$, a continuous
pseudometric on a compact space. For $\sigma,\sigma'\in G^{B_v}$ and $g\in G$
the gauge transformation $g$ at $v$ changes only the links of $B_v$, so by
gauge invariance $h(\sigma',U')=h(g\sigma',U')$ and
\[
|h(\sigma,U)-h(\sigma',U')|\le\Lip(h)\Bigl(\sum_{e\in B_v}\delta(\sigma_e,
(g\sigma')_e)+\delta(U_i,U'_i)\Bigr);
\]
the infimum over $g$ gives $\Lip(h)(\bar d_v(\sigma,\sigma')+\delta(U_i,U'_i))$.
By Lemma~\ref{lem:stars}(iv),~(v), the laws $\gamma_v(\cdot\mid U)$ and
$\gamma_v(\cdot\mid U')$ differ only through the two-path of one plaquette,
which moves by $\delta(U_i,U'_i)$, if $v$ is one of the six vertices of
Lemma~\ref{lem:stars}(v), and coincide otherwise. Lemma~\ref{lem:onestar}
gives
\[
|\mathcal P_vh(U)-\mathcal P_vh(U')|\le\Lip(h)\,\bigl(1+\gamma\,[v\text{ one of
the six}]\bigr)\,\delta(U_i,U'_i).
\]
For $\mathcal K_s=N_s^{-1}\sum_v\mathcal P_v$, summing over $v$ (exactly one
star contains $i$) gives $\Lip(\mathcal K_sh)\le(1-\kappa_0)\Lip(h)$ with
$\kappa_0=(1-6\gamma)/N_s$. Each $\mathcal P_v$ maps $\mathcal A$ to itself, by
gauge covariance of $\gamma_v$, and for a Lipschitz $g\in\mathcal A$,
$\Var_\mu(g)\le\frac12\iint(g(U)-g(U'))^2\,\mu(dU)\mu(dU')\le2\Lip(g)^2|E|^2$.
So $\Var_\mu(\mathcal K_s^th)\le2(1-\kappa_0)^{2t}\Lip(h)^2|E|^2$. On the
mean-zero part of $\mathcal A$ the operator $\mathcal K_s$ is self-adjoint with
$0\le\mathcal K_s\le I$, and $\Var_\mu(\mathcal K_s^th)=\int\lambda^{2t}\,
d\nu_h(\lambda)$ for its spectral measure $\nu_h$ at $h-\mu(h)$; hence $\nu_h$
lives on $[0,1-\kappa_0]$. Lipschitz functions are dense in $L^2(\mu)$, and
the orthogonal projection onto $\mathcal A$, averaging over the compact gauge
group, which acts link by link by isometries, does not increase $\Lip$; so
such $h$ are dense in $\mathcal A$, and the spectral projection of $\mathcal K_s$
on $(1-\kappa_0,1]$ vanishes on the mean-zero part of $\mathcal A$. Thus
$\ip f{(I-\mathcal K_s)f}\ge\kappa_0\Var_\mu(f)$, and $\ip f{(I-\mathcal P_v)f}=
\mu((f-\mathcal P_vf)^2)$.
\end{proof}

\begin{theorem}\label{thm:stars}
Let $G=\SU(2)$ and $|\beta_W|\le\frac9{25}$. Then for every even $L\ge4$ and
every real $f\in\mathcal A$,
\[
\sum_{v\text{ even}}\mu\bigl((f-\mathcal P_vf)^2\bigr)\ge\frac{14}{2957}
\Var_\mu(f),
\]
and for every $n\ge2$ and even $L\ge2n+4$, $\BP(n,\varpi_n)$ holds with
\[
\varpi_n=\frac{14}{2957}\cdot\frac{(n-1)^4}{n^4\lceil(n-1)^4/2\rceil},\qquad
\varpi_2=\frac{14}{16\cdot2957}.
\]
For $0\le\beta_W\le\frac9{25}$ every torus limit is translation invariant and
reflection positive, and its transfer matrix has spectral gap at least
$\varpi_2/(2e^2q_2)=7/(625\cdot2957\,e^2)$.
\end{theorem}

\begin{proof}
Here $c=|\beta|\le\frac9{100}$, and by Lemma~\ref{lem:svalues}
$1-6c\,s(c)\ge\frac{14}{2957}$; the first claim is
Proposition~\ref{prop:starcond}. For a cube $V_x$, an even vertex $v$ with
$v-x\in[1,n-1]^4$ has its eight links based in $x+[0,n)^4$, so $B_v\subset
V_x$. The $\sigma$-algebra of the links outside $V_x$ is contained in that of
the links outside $B_v$, so $\mu((f-\mathcal P_xf)^2)=\mu(f^2)-\mu((\mathcal
P_xf)^2)\ge\mu(f^2)-\mu((\mathcal P_vf)^2)=\mu((f-\mathcal P_vf)^2)$. Each cube
contains at most $N=\lceil(n-1)^4/2\rceil$ such vertices, and each even $v$ is
in this position in exactly $(n-1)^4$ cubes. Averaging over those $v$ in each
cube,
\[
\sum_x\mu((f-\mathcal P_xf)^2)\ge\frac{(n-1)^4}{N}\sum_{v\text{ even}}
\mu((f-\mathcal P_vf)^2)\ge\frac{(n-1)^4}N\cdot\frac{14}{2957}\Var_\mu(f),
\]
which is $\BP(n,\varpi_n)$. The last claim is Corollary~\ref{cor:latticegap}
with $n=2$ and $q_2=\frac{625}{16}$.
\end{proof}

The same proof works whenever $6c\,s(c)<1$. The exact threshold is the root
$c_*$ of $6c\,s(c)=1$, and Lemma~\ref{lem:svalues}(iii) shows
$c_*>\frac9{100}$.\footnote{Solving~\eqref{eq:sevensys} numerically gives
$c_*\approx0.09023$, that is $\beta_W\approx0.3609$. No statement of the paper
uses this value.} At $\beta_W=\frac13$ the star condition holds with
$6c\,s(c)=\frac{13}{15}$, while the single-link argument of
Sec.~\ref{sec:bpstrong} is exactly marginal there. The constant $\varpi_n$ of
Theorem~\ref{thm:stars} decays like $\frac{28}{2957}n^{-4}$, so that theorem is
a cube condition for each fixed $n$, enough for the lattice gap but not for
Theorem~\ref{thm:B}. The lower faces of a cube are the reason, and the next
proposition shows that they stop every certificate that treats their links
one at a time.

\begin{proposition}[the face barrier]\label{prop:facebarrier}
For a lower face $\{\hat v_i=0\}$ of the cube $V_x$, $\hat v=v-x$, let $F_i$ be
the set of links $(b,\mu)\in V_x$ with $\hat b_i=0$ and $\mu\ne i$, and let
$A_F$ be the matrix on $F_i$ with $(A_F)_{ek}=1$ when $e\ne k$ lie in a common
plaquette all of whose links are in $F_i$. Then $\lambda_{\max}(A_F)\ge
12(1-\frac1n)^2$. Consequently, if a certificate for the response of a cube
treats the links of $F_i$ as single sites with influence matrix $C\ge cA_F$ on
$F_i\times F_i$, the inequalities $0\le\varrho\le C\varrho+\mathfrak s$ admit
solutions with $\sum\varrho$ arbitrarily large as soon as
$12c(1-\frac1n)^2\ge1$; as $n\to\infty$ such a certificate cannot pass
$c=\frac1{12}$, that is $|\beta_W|=\frac13$.
\end{proposition}

\begin{proof}
$|F_i|=3n^3$. A plaquette in a plane $\{\mu,\nu\}\not\ni i$ with base point
$y$, $\hat y_i=0$, has all four links in $F_i$ exactly when its three base
points lie in $x+[0,n)^4$, that is $\hat y_\mu,\hat y_\nu\le n-2$; there are
$3n(n-1)^2$ of them. Two links share at most one plaquette, so
$\mathbf 1^{\mathsf T}A_F\mathbf 1=12\cdot3n(n-1)^2$, and the Rayleigh quotient
of $\mathbf 1$ is $12(1-\frac1n)^2$. If $c\lambda_{\max}(A_F)\ge1$ and $w\ge0$
is a Perron vector of $A_F$, extended by $0$, then $\varrho=tw$ satisfies
$C\varrho\ge tcA_Fw\ge\varrho$ on $F_i$ for every $t>0$.
\end{proof}

For single-link updates the hypothesis $C\ge cA_F$ holds with $c=|\beta|$: at
$b=0$ the derivative of $b\mapsto\nu_b$ has Kantorovich norm exactly
$\frac14$ (Lemma~\ref{q-lem:rate}(i)), and the staples of a face link can
cancel while one of them moves. The forest of Sec.~\ref{sec:bpstrong} does not
help, since its links $T_x$ have interior heads and the face links are free.
What does help is to let the gauge transformation act at the vertices of the
faces as well, and to pay for it.

\subsection{The cube as a block system}

Fix $x$ and $n\ge2$, and write $\hat v=v-x$, computed in a window
$x+[-2,n+2]^4$ that embeds in $\Z^4$ because $L\ge2n+4$. Let $V=V_x$ and
$Q=\{v:\hat v\in[0,n-1]^4\}$. Call $v$ \emph{lower} if $v\in Q$, with zero set
$Z(v)=\{a:\hat v_a=0\}$ and $z(v)=|Z(v)|$, and \emph{upper} if $\hat v\in
[0,n]^4$ has exactly one coordinate equal to $n$. If $k$ is a link with odd end
point $o$ and direction $\mu$, its six plaquettes are indexed by $(\nu,
\epsilon)$, $\nu\ne\mu$, $\epsilon=\pm1$; the even corner of the plaquette
$(\nu,\epsilon)$ other than the even end point of $k$ is $o+\epsilon e_\nu$,
and there $k$ is a factor of the two-path of this plaquette only
(Lemma~\ref{lem:stars}(v)).

\begin{lemma}\label{lem:cubestars}
Let $n\ge2$.
\begin{enumerate}[label=\textup{(\roman*)},leftmargin=*]
\item If $v$ is lower and even, $B_v\cap V$ consists of $B_v$ minus the
$z(v)$ \emph{outside links} $(v-e_a,a)$, $a\in Z(v)$; we call the directions
$-a$, $a\in Z(v)$, \emph{fixed} and the others \emph{free}. If $v$ is upper
and even with $\hat v_a=n$, then $B_v\cap V=\{(v-e_a,a)\}$. Other even vertices
carry no link of $V$. Hence $V$ is partitioned into the \emph{blocks}
$B_v\cap V$, $v$ even and lower or upper, and the \emph{owner} $w(k)$ of
$k\in V$ is its even end point.
\item Let $k\in V$ have odd end point $o$, direction $\mu$ and owner
$o+\varsigma e_\mu$, $\varsigma=\pm1$, and let $p$ be its plaquette
$(\nu,\epsilon)$, $v=o+\epsilon e_\nu$. The links of $p$ at $v$ are in the
directions $-\epsilon\nu$ and $\varsigma\mu$. If $v$ is lower, both lie in
$V$. If $v$ is upper with $\hat v_a=n$, the links of $p$ at $v$ other than
$(v-e_a,a)$ lie outside $V$.
\item Let $v$ be lower and even and $a\in Z(v)$. The outside link $(v-e_a,a)$
is a factor of no two-path $h_p$ at a lower or upper vertex. If a plaquette
$p$ at $v$ has one of its links at $v$ outside $V$, then both links of $p$ at
the opposite even corner lie outside $V$.
\end{enumerate}
\end{lemma}

\begin{proof}
(i) The outgoing links $(v,b)$ of a lower $v$ have base point in $Q$, and an
incoming link $(v-e_b,b)$ lies in $V$ if and only if $\hat v_b\ge1$. If
$\hat v_a=n$ for exactly one $a$, no outgoing link and no incoming link
$(v-e_b,b)$, $b\ne a$, has its base point in $Q$, while $(v-e_a,a)$ has. If
$\hat v\notin[0,n]^4$ or two coordinates equal $n$, no link at $v$ has its base
point in $Q$.

(ii) The corners of $p$ are $o$, $o+\varsigma e_\mu$, $o+\varsigma e_\mu+
\epsilon e_\nu$ and $v$, so the links at $v$ go to $o$, in direction
$-\epsilon\nu$, and to $v+\varsigma e_\mu$, in direction $\varsigma\mu$. Let $v$
be lower. If $\epsilon=+1$ the first link is $(o,\nu)$, and $\hat o=\hat v-e_\nu$
lies in $Q$, its $\nu$-coordinate being $\hat o_\nu\ge0$; if $\epsilon=-1$ it is
$(v,\nu)$, with base point $v\in Q$. The second link is $(v,\mu)$ if
$\varsigma=+1$; if $\varsigma=-1$ it is $(v-e_\mu,\mu)$, which lies in $V$
because $\hat v_\mu=\hat o_\mu\ge1$, as $k=(o-e_\mu,\mu)\in V$. If $v$ is upper,
(i) gives the claim.

(iii) The plaquettes of $(v-e_a,a)$ have the even corners $v$ and
$v-e_a\pm e_\nu$, whose $a$-coordinate is $-1$, so these are neither lower nor
upper; and a link is a factor of a two-path only at the even corner opposite
to its own even end point. If $p$ at $v$ is spanned by $-a$ and $\tau$, the
opposite corner $v'=v-e_a+\tau$ and both its neighbours in $p$, $v-e_a$ and
$v+\tau$, are joined to it by links based at $a$-coordinate $-1$.
\end{proof}

For a block vertex $v$ the conditional law $\gamma_v(\cdot\mid\sigma,U)$ of
the links of $B_v\cap V$, given the other links of $V$ (values $\sigma$) and
the links off $V$ (values $U$), has density proportional to the product of
the weights $e^{4\beta\ip{h_p}{X_\sigma^{-1}X_\tau}}$ of the plaquettes $p$ at
$v$ with a link at $v$ in $V$. It depends on $(\sigma,U)$ only through these
two-paths and the values of the links at $v$ off $V$.

\subsection{The soft quotient}

At a vertex of a lower face the gauge transformation moves the outside links
at that vertex, and the conditional law of the cube does not see it. The
following pseudometric allows it anyway, at the price of the displacement it
causes on those links. The price is legitimate in the Lipschitz contraction
argument of Proposition~\ref{prop:interiorbp}, because the test function
depends on all links.

\begin{definition}\label{def:dtilde}
For a lower even $v$ and configurations $X,X'$ of $B_v\cap V$, written as
outward transporters, put
\[
\tilde d_v(X,X')=\min_{g\in G}\Bigl[\sum_{\sigma\text{ free}}\delta(gX_\sigma,
X'_\sigma)+z(v)\,\delta(g,1)\Bigr],
\]
and for an upper even $v$ let $\tilde d_v$ be $\delta$ on its single link. For
$\sigma,\sigma'\in G^V$ let $D(\sigma,\sigma')=\sum_v\tilde d_v(\sigma_{B_v},
\sigma'_{B_v})$, the sum over the block vertices, and for $j\notin V$
\[
\tilde R_x(j)=\sup\Bigl\{\frac{W^D(\gamma^0_{V,U},\gamma^0_{V,U'})}
{\delta(U_j,U'_j)}:U,U'\text{ differ only at }j\Bigr\},
\]
where $\gamma^0_{V,U}=\gamma^0_V(\cdot\mid U)$ and $W^D$ is the Kantorovich
distance for the cost $D$.
\end{definition}

The minimum exists by compactness, and $\tilde d_v$ and $D$ are continuous,
symmetric, since $\delta(gX,X')=\delta(X,g^{-1}X')$ and
$\delta(g,1)=\delta(g^{-1},1)$, and satisfy the triangle inequality, by
composing the minimising elements. For $z(v)=0$, $\tilde d_v=\bar d_v$.

\begin{lemma}\label{lem:softlip}
Let $h$ be a function on $G^E$, gauge invariant at every configuration, with
$|h(W)-h(W')|\le\Lip(h)\sum_e\delta(W_e,W'_e)$. If $U,U'$ differ only at
$j\notin V$, then for all $\sigma,\sigma'\in G^V$
\[
|h(\sigma,U)-h(\sigma',U')|\le\Lip(h)\bigl(D(\sigma,\sigma')+\delta(U_j,U'_j)
\bigr).
\]
\end{lemma}

\begin{proof}
For each lower even $v$ choose $g_v$ attaining $\tilde d_v(\sigma_{B_v},
\sigma'_{B_v})$, and let $g$ be the gauge transformation equal to $g_v$ at
these vertices and to $1$ elsewhere; then $h(\sigma,U)=h(g\cdot(\sigma,U))$. A
link of $V$ has exactly one even end point $v$, and $g$ acts on it by $X\mapsto
g_vX$ on its outward transporter at $v$, or trivially if $v$ is upper. A link
$e\notin V$ has at most one end point where $g\ne1$, namely an even lower
$v$, and then $e$ is one of the $z(v)$ outside links at $v$, by
Lemma~\ref{lem:cubestars}(i); there $\delta((g\cdot U)_e,U'_e)\le\delta(g_v,1)+
\delta(U_e,U'_e)$ by bi-invariance. Summing over all links,
$|h(g\cdot(\sigma,U))-h(\sigma',U')|\le\Lip(h)\bigl(\sum_v\tilde d_v+
\delta(U_j,U'_j)\bigr)$.
\end{proof}

\subsection{Block responses}

\begin{lemma}[Dobrushin on $S^3$]\label{lem:dobS3}
Assume~\eqref{eq:Q}. Let $S$ be finite and $\nu,\nu'$ probability measures on
$G^S$ with continuous positive densities with respect to Haar measure, such
that for every $s\in S$ the conditional law of $x_s$ given $x_{S\setminus s}$
is $\nu_{B_s(x)}$ under $\nu$ and $\nu_{B'_s(x)}$ under $\nu'$, where
$|B_s(x)-B_s(y)|\le4|\beta|\sum_ta_{st}\delta(x_t,y_t)$ and
$|B_s(x)-B'_s(x)|\le4|\beta|\eta_s$ for all $x,y$, with $a_{st}\ge0$,
$a_{ss}=0$ and $\eta_s\ge0$. Then there is a coupling $\pi$ of $\nu,\nu'$ whose
$\varrho_s=\int\delta(x_s,x'_s)\,d\pi$ satisfy $\varrho\le cA\varrho+c\eta$
componentwise. Consequently:
\begin{enumerate}[label=\textup{(\alph*)},leftmargin=*]
\item if $\sum_sa_{st}\le d$ for every $t$ and $cd<1$, then
$\sum_s\varrho_s\le c\sum_s\eta_s/(1-cd)$;
\item if the spectral radius of $cA$ is less than $1$ and $y\in[0,\infty)^S$
satisfies $y\ge cAy+c\eta$, then $\varrho\le y$;
\item if $y\ge0$ satisfies $y\ge c_0Ay+\eta$ and $c\le c_0$, then $cy\ge
cA(cy)+c\eta$, so $\varrho\le cy$ when the spectral radius of $cA$ is less
than $1$.
\end{enumerate}
\end{lemma}

\begin{proof}
Let $\pi$ minimise $\int\sum_s\delta(x_s,x'_s)\,d\pi$ over couplings. Fix $s$
and replace under $\pi$ the pair $(x_s,x'_s)$ by a sample of an optimal
coupling of $\nu_{B_s(x)}$ and $\nu_{B'_s(x')}$, chosen measurably in
$(x_{S\setminus s},x'_{S\setminus s})$~\cite[Cor.~5.22]{Vil09}. This is again
a coupling of $\nu,\nu'$, and only the $s$-term of the cost changes, so by
minimality,~\eqref{eq:Q} and the triangle inequality, $\varrho_s\le\theta\int
|B_s(x)-B'_s(x')|\,d\pi\le c\sum_ta_{st}\varrho_t+c\eta_s$. (a) Summing over
$s$, $\sum\varrho\le cd\sum\varrho+c\sum\eta$. (b) The vector $y-\varrho$
satisfies $y-\varrho\ge cA(y-\varrho)$, hence $y-\varrho\ge(cA)^K(y-\varrho)
\to0$. (c) $cy-c^2Ay\ge c(y-c_0Ay)\ge c\eta$; apply (b).
\end{proof}

From now on $c\le c_0=\frac7{80}$, that is $|\beta_W|\le\frac7{20}$, and
$\kappa_0=s(c_0)=\frac{3480}{1933}<1.8004$, so that $6c_0\kappa_0=
\frac{1827}{1933}$. For a direction $a$ and a family $h$ of two-paths let
$\mu^{(a)}_h$ be the law on $G^{S_a}$ with density proportional to
$\prod_{\{\sigma,\tau\}}e^{4\beta\ip{h_{\sigma\tau}}{Z_\sigma^{-1}Z_\tau}}$,
$Z_a=1$, as in the proof of Lemma~\ref{lem:onestar}. By that proof,
Lemma~\ref{lem:svalues}(ii) and Lemma~\ref{lem:dobS3}(c), if $h,h'$ differ only
at the pair $\{a,b\}$ there is a coupling of $\mu^{(a)}_h$ and $\mu^{(a)}_{h'}$
with
\begin{equation}\label{eq:sevenbound}
\int\sum_{\sigma\in S_a}\delta(Z_\sigma,Z'_\sigma)\le\kappa_0c\,
|h_{ab}-h'_{ab}|\qquad(c\le c_0).
\end{equation}

\begin{lemma}\label{lem:partial}
Let $c\le c_0$ and $v$ a block vertex, and write $\gamma^{h,Y}$ for its block
law as a function of the two-paths $h$ at $v$ and of the values $Y$ of the
links at $v$ off $V$. Put $\kappa_1=\kappa_0$, $\kappa_2=\frac94$,
$\kappa_3=\frac{40}{13}$ and $\kappa_4=\frac{160}{59}$.
\begin{enumerate}[label=\textup{(\alph*)},leftmargin=*]
\item Let $v$ be lower, and let $(h,Y)$ and $(h',Y)$ differ only at the
two-path $h_p$ of a plaquette $p$ whose two links at $v$ are free. Then
$W^{\tilde d_v}(\gamma^{h,Y},\gamma^{h',Y})\le\kappa_{z(v)}c\,|h_p-h'_p|$, with
$\kappa_0$ for $z(v)=0$.
\item In particular, for $z(v)=1$ the constant is the full-star constant
$\kappa_0c$.
\item Let $v$ be lower and $p$ a plaquette at $v$ with one link at $v$ free and
one fixed. Then the bound of \textup{(a)} holds with $\kappa_0c$ if $z(v)=1$
and with $\frac{16}9c$ if $z(v)\ge2$.
\item Let $v$ be lower and let $(h,Y)$ and $(h,Y')$ differ only at one
outside link, by $\eta$ in $\delta$. Then $W^{\tilde d_v}(\gamma^{h,Y},
\gamma^{h,Y'})\le\eta$.
\item Let $v$ be upper. If $(h,Y)$ and $(h',Y')$ differ only at one two-path
$h_p$ with $p\ni(v-e_a,a)$, or only at one outside link at $v$, the distance is
at most $c|h_p-h'_p|$, respectively $c\eta$.
\end{enumerate}
\end{lemma}

\begin{proof}
\emph{Representation.} Let $v$ be lower and $F$ its set of fixed directions,
let $\Phi_h(X)=\prod e^{4\beta\ip{h_{\sigma\tau}}{X_\sigma^{-1}X_\tau}}$ over all
$24$ adjacent pairs, so that $\Phi_h(gX)=\Phi_h(X)$, and note that
$\gamma^{h,Y}$ has density proportional to $\Phi_h$ evaluated at $X_F=Y_F$. Fix
a free direction $\sigma_0$. If $F=\emptyset$, the map
$X\mapsto(X_{\sigma_0},(X_{\sigma_0}^{-1}X_\rho)_{\rho\ne\sigma_0})$ preserves
Haar measure and carries $\gamma^h$ to $\mathrm{Haar}\otimes\mu^{(\sigma_0)}_h$,
as in the proof of Lemma~\ref{lem:onestar}. If $F\ne\emptyset$, fix $f_1\in F$
and let $\Psi(X)$ consist of $W_\rho=X_{\sigma_0}^{-1}X_\rho$ for the free
$\rho\ne\sigma_0$ and of $W_m=X_{\sigma_0}^{-1}Y_{f_1}$. Then $\Psi$ is a
bijection of the free coordinates, with inverse $X_{\sigma_0}=Y_{f_1}W_m^{-1}$,
$X_\rho=Y_{f_1}W_m^{-1}W_\rho$, and it preserves Haar measure: $X_{\sigma_0}
\mapsto X_{\sigma_0}^{-1}Y_{f_1}$ does, and for fixed $X_{\sigma_0}$ so does
each $X_\rho\mapsto X_{\sigma_0}^{-1}X_\rho$. Since $\Phi_h(X)=\Phi_h(
X_{\sigma_0}^{-1}X)$ and $X_{\sigma_0}^{-1}X_f=W_mK_f$ with $K_f=Y_{f_1}^{-1}
Y_f$, the image of $\gamma^{h,Y}$ has density proportional to $\Phi_h$ at
$X_{\sigma_0}=1$, $X_\rho=W_\rho$, $X_f=W_mK_f$. If $(W,W')$ couples the images
for $(h,Y)$ and $(h',Y)$ and $X=\Psi^{-1}W$, $X'=\Psi^{-1}W'$, then
$g=Y_{f_1}W'^{-1}_mW_mY_{f_1}^{-1}$ satisfies $gX_{\sigma_0}=X'_{\sigma_0}$,
$\delta(gX_\rho,X'_\rho)=\delta(W_\rho,W'_\rho)$ and $\delta(g,1)=
\delta(W_m,W'_m)$, so
\begin{equation}\label{eq:transfer}
\tilde d_v(X,X')\le\sum_\rho\delta(W_\rho,W'_\rho)+z(v)\,\delta(W_m,W'_m).
\end{equation}

(a), (b) for $z(v)\le1$. Take $\sigma_0$ and $\tau$ to be the two directions of
$p$ at $v$. For $z(v)=1$, $F=\{f_1\}$, the image is exactly
$\mu^{(\sigma_0)}_h$, the coordinate $W_m$ being $Z_{f_1}$, and the change of
$h_p=h_{\sigma_0\tau}$ is a change of the field on $Z_\tau$;
\eqref{eq:sevenbound} and~\eqref{eq:transfer} give $\kappa_0c|h_p-h'_p|$. For
$z(v)=0$ couple with the same $X_{\sigma_0}$ and use $g=1$; this is
Lemma~\ref{lem:onestar} with~\eqref{eq:sevenbound}.

(a) for $z(v)=2$. Let $F=\{-i,-j\}$; the free directions are $P=+i$, $Q=+j$,
$K_\pm=\pm k$ and $L_\pm=\pm l$. Up to the symmetries $i\leftrightarrow j$,
$k\leftrightarrow l$, $k\mapsto-k$ and $l\mapsto-l$ of the star, which fix
$F$, the pair of $p$ is one of $\{P,Q\}$, $\{P,K_+\}$ and $\{K_+,L_+\}$. In the
image the terms of the density are: for adjacent free $\rho,\rho'\ne\sigma_0$,
$\ip h{W_\rho^{-1}W_{\rho'}}$; for $\{\sigma_0,\rho\}$, a field on $W_\rho$; for
$\{f,\rho\}$, $\ip h{(W_mK_f)^{-1}W_\rho}=\ip{K_fh}{W_m^{-1}W_\rho}$, linear in
$W_m$ and in $W_\rho$; for $\{\sigma_0,f\}$, a field on $W_m$; and for
$\{-i,-j\}$, a constant. So Lemma~\ref{lem:dobS3} applies with
$a_{\rho\rho'}=1$ for adjacent free $\rho,\rho'\ne\sigma_0$,
$a_{\rho m}=a_{m\rho}$ equal to the number of $f\in F$ adjacent to $\rho$
(one for $P,Q$, two for $K_\pm,L_\pm$), and $\eta=|h_p-h'_p|e_\tau$. Every row
sum of $A$ is at most $9$, so $c_0A$ has spectral radius at most
$\frac{63}{80}$. By~\eqref{eq:transfer} the cost is at most
$\sum_\rho\varrho_\rho+2\varrho_m$. We give $\hat y\ge c_0A\hat y+e_\tau$ and
use Lemma~\ref{lem:dobS3}(c); each item lists $\hat y$ and the inequalities
multiplied by $80$ and by the denominator of $\hat y$.

\emph{The pair $\{P,Q\}$}, $\sigma_0=P$, $\tau=Q$. The spins are $Q$, $K_\pm$,
$L_\pm$, each of the last four adjacent to $Q$ and to the two of the other
axis, and $m$; $100\hat y=(Q,K_\pm,L_\pm,m)=(108,16,21)$. The inequalities
are $8640\ge7(4\cdot16+21)+8000=8595$, $1280\ge7(108+2\cdot16+2\cdot21)=1274$
and $1680\ge7(108+8\cdot16)=1652$, and the cost is at most
$(108+4\cdot16+2\cdot21)/100=2.14$.

\emph{The pair $\{P,K_+\}$}, $\sigma_0=K_+$, $\tau=P$. The spins are $P$ and
$Q$, adjacent to each other and to $K_-$, $L_\pm$; $K_-$, adjacent to $P,Q,
L_\pm$; $L_\pm$, adjacent to $P,Q,K_-$; and $m$; $200\hat y=(P,Q,K_-,L_\pm,m)=
(216,32,35,32,40)$. The inequalities are $17280\ge7(32+35+64+40)+16000=17197$,
$2560\ge7(216+35+64+40)=2485$, $2800\ge7(216+32+64+80)=2744$,
$2560\ge7(216+32+35+80)=2541$ and $3200\ge7(216+32+70+128)=3122$, and the
cost is at most $(216+32+35+64+80)/200=2.135$.

\emph{The pair $\{K_+,L_+\}$}, $\sigma_0=K_+$, $\tau=L_+$, with the same spins
and adjacencies; $200\hat y=(P,Q,K_-,L_+,L_-,m)=(32,32,36,219,19,54)$. The
inequalities are $2560\ge7(32+36+219+19+54)=2520$ for $P$ and for $Q$,
$2880\ge7(64+219+19+108)=2870$, $17520\ge7(64+36+108)+16000=17456$,
$1520\ge7(64+36+108)=1456$ and $4320\ge7(64+72+2\cdot238)=4284$, and the cost
is at most $(64+36+219+19+108)/200=2.23$.

In all cases the response is at most $2.23\,c|h_p-h'_p|\le\kappa_2c|h_p-h'_p|$.

(a) for $z(v)\ge3$, and (c), (d) for $z(v)\ge2$. Use the coordinates
$X_\sigma$, $\sigma$ free, directly, and $g=1$ in $\tilde d_v$. Each conditional
law is $\nu_B$ with $B=4\beta$ times the sum of the unit vectors of the
plaquettes at $v$ through $\sigma$, those through free neighbours depending
isometrically on them; so Lemma~\ref{lem:dobS3}(a) applies with $A$ the
adjacency matrix of the free directions, whose largest degree is $d=5,4,3$
for $z(v)=2,3,4$ (for $z=2$, $+i$ is adjacent to $+j$, $\pm k$, $\pm l$; for
$z=3$, $+i$ to $+j$, $+k$, $\pm l$; for $z=4$ the free directions are the four
positive ones). A change of $h_p$ with both links free gives $\sum\eta=
2|h_p-h'_p|$, with one link free $\sum\eta=|h_p-h'_p|$, and a change of the
outside link in direction $f$ by $\eta$ gives $\eta$ on each free neighbour of
$f$, at most $d$ of them. With $c\le c_0$: $\frac{2c}{1-4c}\le\frac{40}{13}c$,
$\frac{2c}{1-3c}\le\frac{160}{59}c$, $\frac c{1-cd}\le\frac c{1-5c_0}=
\frac{16}9c$ and $\frac{dc}{1-dc}\le\frac{5c_0}{1-5c_0}=\frac79$.

(c), (d) for $z(v)=1$. Let $F=\{f\}$ and $X_f=Y$. The map $X\mapsto Z=(Y^{-1}
X_\rho)_\rho$ carries $\gamma^{h,Y}$ to $\mu^{(f)}_h$, which does not depend on
$Y$. A change of $h_p$, $p$ spanned by $f$ and $\tau$, is a change of the
field on $Z_\tau$, so~\eqref{eq:sevenbound} gives a coupling with
$\int\sum\delta(YZ_\rho,YZ'_\rho)\le\kappa_0c|h_p-h'_p|$, with $g=1$. For a
change $Y\to Y'$ couple $X=YZ$ and $X'=Y'Z$ and take $g=Y'Y^{-1}$; then
$\tilde d_v(X,X')\le\delta(Y'Y^{-1},1)=\delta(Y,Y')$.

(e) The block law is $\nu_b$ with $b=4\beta$ times the sum of the six vectors
of the plaquettes through $(v-e_a,a)$; such a change moves $b$ by at most
$4|\beta||h_p-h'_p|$, respectively $4|\beta|\eta$, since an outside link lies
in at most one plaquette with $(v-e_a,a)$. Apply~\eqref{eq:Q}.
\end{proof}

Part (b) is the point of the soft quotient: a star with one fixed link is
exactly as good as a full star, and the price $z(v)\delta(g,1)$ never feeds
back into the system, because an outside link of a lower vertex is a factor
of no two-path at a block (Lemma~\ref{lem:cubestars}(iii)). The margin is lost
at the edges of the cube, $z(v)=2$, and the next lemma pays for them with
weights.

\subsection{The column inequality}

For a block vertex $v$ and a plaquette $p$ at $v$ let $\Gamma_v(p)$ be the
constant of Lemma~\ref{lem:partial}(a), (c) or (e), and $\Gamma_v(p)=0$ if no
link of $p$ at $v$ lies in $V$. Put $u_v=1$ for upper $v$ and $u_v=u_{z(v)}$
for lower $v$, with
\[
u_0=u_1=1,\quad u_2=\tfrac56,\quad u_3=\tfrac7{10},\quad u_4=\tfrac35;
\qquad\lambda=\tfrac{97}{100}.
\]

\begin{lemma}\label{lem:column}
Let $n\ge4$ and $c\le c_0$. For every $k\in V$, with plaquettes
$p_{(\nu,\epsilon)}$ and even corners $v_{(\nu,\epsilon)}=o+\epsilon e_\nu$ as
above,
\[
\sum_{(\nu,\epsilon)}u_{v_{(\nu,\epsilon)}}\Gamma_{v_{(\nu,\epsilon)}}
(p_{(\nu,\epsilon)})\le\lambda\,u_{w(k)} .
\]
\end{lemma}

\begin{proof}
Put $f(z)=u_z\kappa_zc$: $f(0)=f(1)=\kappa_0c<1.8004\,c$, $f(2)=\frac{15}8c$,
$f(3)=\frac{28}{13}c<2.1539\,c$ and $f(4)=\frac{96}{59}c<1.6272\,c$, all at
least $c$. By Lemma~\ref{lem:cubestars}(ii) a lower $v_{(\nu,\epsilon)}$ has
both links of $p$ free, so it contributes $f(z(v_{(\nu,\epsilon)}))$; an upper
one contributes at most $c$, and other vertices contribute $0$. Let $o$, $\mu$
and $\varsigma$ be as in Lemma~\ref{lem:cubestars}(ii).

\emph{Case A: $\hat o$ has a coordinate equal to $n$.} Since $k\in V$ is
$(o,\mu)$, which needs $o\in Q$, or $(o-e_\mu,\mu)$, which needs $\hat o-e_\mu
\in[0,n-1]^4$, this coordinate is $\hat o_\mu$, and $w(k)=o-e_\mu$ with
$\hat w_\mu=n-1\ge1$; so $z(w)\le3$ and $u_w\ge\frac7{10}$. Every
$v_{(\nu,\epsilon)}$ has $\mu$-coordinate $n$, so the sum is at most
$6c\le6c_0=0.525<\frac7{10}\cdot0.97$.

\emph{Case B: $o\in Q$.} Let $Z=z(o)$ be the number of zeros of $\hat o$,
$a=\#\{\nu\ne\mu:\hat o_\nu=0\}$ and $b=\#\{\nu\ne\mu:\hat o_\nu=1\}$. As
$n\ge4$, the classes $\{0\}$, $\{1\}$, $[2,n-2]$ and $\{n-1\}$ are disjoint.
For $\nu\ne\mu$: if $\hat o_\nu=0$, $\epsilon=-1$ gives no block and
$\epsilon=+1$ a lower vertex with $z=Z-1$; if $\hat o_\nu=1$, $\epsilon=+1$
gives $z=Z$ and $\epsilon=-1$ gives $z=Z+1$; if $2\le\hat o_\nu\le n-2$ both
give $z=Z$; and if $\hat o_\nu=n-1$, $\epsilon=-1$ gives $z=Z$ and
$\epsilon=+1$ an upper vertex, which contributes at most $c\le f(Z)$. Hence the
sum is at most
\[
a\,f(Z-1)+b\bigl(f(Z)+f(Z+1)\bigr)+2(3-a-b)f(Z).
\]
If $\hat o_\mu=0$, then $\varsigma=+1$, $Z=a+1$ and $z(w)=a$, and the sum is at
most $h(a,b)=af(a)+b(f(a+1)+f(a+2))+2(3-a-b)f(a+1)$, to be compared with
$\lambda u_a$. If $\hat o_\mu\ge1$, then $Z=a$ and $w$ is upper, with $u_w=1$,
or lower with $z(w)\in\{a,a+1\}$, so $u_w\ge u_{a+1}$; the sum is at most
$g(a,b)=af(a-1)+b(f(a)+f(a+1))+2(3-a-b)f(a)$, to be compared with $\lambda
u_{a+1}$. Here $a+b\le3$. Since $f(0)=f(1)\le f(2)\le f(3)$, the quantities
$h(0,b)$, $g(1,b)$, $h(1,b)$ and $g(2,b)$ increase with $b$, and
$g(0,b)=6f(0)$. Every term is $c$ times a constant, so it suffices to take
$c=c_0$; with $\hat\lambda=\lambda/c_0=\frac{388}{35}>11.0857$ and $f$ in units
of $c$, the cases are
\begin{align*}
h(0,3)&=3f(0)+3f(2)<5.4012+5.625=11.0262<\hat\lambda,\\
g(0,b)&=6f(0)<10.8024<\hat\lambda,\\
g(1,2)&=3f(0)+2f(2)<9.1512<\tfrac56\hat\lambda,\quad\tfrac56\hat\lambda>9.2380,\\
h(1,2)&=f(0)+2f(2)+2f(3)<9.8582<\hat\lambda,\\
g(2,1)&=2f(0)+f(2)+f(3)<7.6297<\tfrac7{10}\hat\lambda,\quad
\tfrac7{10}\hat\lambda>7.7599,\\
h(2,0)&=2f(2)+2f(3)<8.0578<\tfrac56\hat\lambda,\\
h(2,1)&=2f(2)+f(3)+f(4)<7.5311<\tfrac56\hat\lambda,\\
g(3,0)&=3f(2)=5.625<\tfrac35\hat\lambda,\quad\tfrac35\hat\lambda>6.6514,\\
h(3,0)&=3f(3)<6.4617<\tfrac7{10}\hat\lambda.
\end{align*}
All inequalities hold.
\end{proof}

\subsection{The cube condition uniformly in the cube size}

\begin{proposition}\label{prop:starresponse}
Let $c\le c_0$, $n\ge4$ and $L\ge2n+4$ even. For every link $j$,
\[
\sum_{x:\,j\notin V_x}\tilde R_x(j)\le259\,n^3 .
\]
\end{proposition}

\begin{proof}
Fix $x$, $j\notin V=V_x$ and $U,U'$ differing only at $j$, and write
$\eta=\delta(U_j,U'_j)$. Let $\pi$ minimise $\int D\,d\pi$ over couplings of
$\gamma^0_{V,U}$ and $\gamma^0_{V,U'}$, and put $\varrho_v=\int\tilde d_v\,d\pi$.
As in Lemma~\ref{lem:dobS3}, replacing the block of $v$ by an optimal
coupling for $\tilde d_v$ of $\gamma_v(\cdot\mid\sigma,U)$ and
$\gamma_v(\cdot\mid\sigma',U')$ changes only the $v$-term, so $\varrho_v\le\int
\Omega_v\,d\pi$ with $\Omega_v=W^{\tilde d_v}(\gamma_v(\cdot\mid\sigma,U),
\gamma_v(\cdot\mid\sigma',U'))$.

\emph{Pointwise bound.} Changing the two-paths at $v$ one at a time, and then
the outside link $j$ if it is a link at $v$, the triangle inequality for
$W^{\tilde d_v}$ and Lemma~\ref{lem:partial} give $\Omega_v\le\sum_p\Gamma_v(p)
|h_p(\sigma,U)-h_p(\sigma',U')|+r_v\eta$, where $r_v\le1$ if $v$ is lower and
$j$ is an outside link at $v$, $r_v\le c$ if $v$ is upper and $j$ is a link at
$v$, and $r_v=0$ otherwise. For each block vertex $w$ fix $g^*_w$ attaining
$\tilde d_w(\sigma_{B_w},\sigma'_{B_w})$, with $g^*_w=1$ for upper $w$, and put
$\delta_k=\delta(g^*_{w(k)}X_k(\sigma),X_k(\sigma'))$ for $k\in V$, so that
$\sum_{k\in B_w\cap V}\delta_k\le\tilde d_w$. If both links of $p$ at the
opposite corner $v'$ lie in $V$, then $v'$ is lower, an upper vertex carrying
one link of $V$, and by invariance of $h_p$ under $X\mapsto g^*_{v'}X$ at $v'$,
$|h_p(\sigma,U)-h_p(\sigma',U')|$ is at most the sum of their $\delta_k$. If
exactly one lies in $V$, the other is off $V$ and is not an outside link of a
lower vertex (Lemma~\ref{lem:cubestars}(iii)), so $v'$ is upper and the bound
is $\delta_k+[j\in p]\eta$. If none lies in $V$, the bound is $[j\in p]\eta$.
Each $k\in V$ is a factor of a two-path at $v$ for at most one $p$ at $v$, and
these $v$ are the vertices $v_{(\nu,\epsilon)}$ of $k$, so
Lemma~\ref{lem:column} gives, pointwise,
\[
\sum_vu_v\Omega_v\le\lambda\sum_{k\in V}u_{w(k)}\delta_k+B\eta\le\lambda\sum_w
u_w\tilde d_w+B\eta,
\]
with $B=\sum_{p\ni j}u_{v(p)}\Gamma_{v(p)}(p)+\max_vu_vr_v$, where $v(p)$ is
the even corner of $p$ opposite to the even end point of $j$.

\emph{Constants.} By Lemma~\ref{lem:partial}, $u_v\Gamma_v(p)\le\frac{28}{13}c$
in all cases ($\max_zf(z)=f(3)$; with one fixed link $\kappa_0c$ or
$\frac{16}9c$; upper vertices $c$), and only the plaquettes $p\ni j$
containing a link of $V$ contribute. Integrating, $\sum_vu_v\varrho_v\le
\lambda\sum_vu_v\varrho_v+B_x(j)\eta$, with $B_x(j)=\frac{28}{13}c\,P_x(j)+
R_x(j)$, where $P_x(j)$ is the number of plaquettes $p\ni j$ with a link in
$V_x$, and $R_x(j)=1$ if $j$ is an outside link at a lower vertex of the
cube, $c$ if $j$ is a link off $V_x$ at an upper even vertex of the cube, and
$0$ otherwise. All terms are finite, so $\sum_vu_v\varrho_v\le
B_x(j)\eta/(1-\lambda)$ and
\[
W^D(\gamma^0_{V,U},\gamma^0_{V,U'})\le\sum_v\varrho_v\le\frac{B_x(j)\,\eta}
{u_4(1-\lambda)}=\frac{500}9B_x(j)\,\eta .
\]

\emph{Summation over $x$.} A plaquette counted in $P_x(j)$ contains a partner
$k$ of $j$ with $k\in V_x$, and distinct plaquettes give distinct partners.
For $j=(b,\nu)$ the eighteen partners have base points $b+\Delta$ with
$\sum\|\Delta\|_1=18$, $\Delta\in\{0,e_\mu,e_\nu,-e_\mu,-e_\mu,e_\nu-e_\mu\}$
for each $\mu\ne\nu$, and $\#\{x:k\in V_x,\ j\notin V_x\}=|Q_0\setminus(Q_0-
\Delta)|\le n^3\|\Delta\|_1$ with $Q_0=[0,n)^4$, by moving one coordinate at a
time; so $\sum_xP_x(j)\le18n^3$. For $R_x$, let $v$ be the even end point of
$j$. If $v$ is the head of $j=(v-e_a,a)$, it is a lower vertex with $a\in
Z(v)$ for exactly $n^3$ cubes, and an upper vertex of the cube with $j$ off
$V_x$ for at most $3n^3$ cubes, the coordinate equal to $n$ not being $a$; if
$v$ is the base point of $j$, it is never lower with $j\notin V_x$ and is upper
for at most $4n^3$ cubes. So $\sum_xR_x(j)\le(1+3c)n^3$. With $c\le c_0$,
\[
\sum_x\tilde R_x(j)\le\frac{500}9\Bigl(\frac{28}{13}\cdot18c_0+1+3c_0\Bigr)n^3
=\frac{500}9\cdot\frac{4841}{1040}\,n^3<258.61\,n^3 .\qedhere
\]
\end{proof}

\begin{theorem}\label{thm:starsuniform}
Let $G=\SU(2)$ and $|\beta_W|\le\frac7{20}$.
\begin{enumerate}[label=\textup{(\roman*)},leftmargin=*]
\item For every $n\ge260$ and even $L\ge2n+4$, $\BP(n,1-259/n)$ holds; in
particular $\BP(n,\frac12)$ holds for $n\ge518$.
\item For every $n\ge2$ and even $L\ge2n+4$, $\BP(n,\varpi)$ holds with
$\varpi=\frac{106}{1933}\cdot517^{-4}$, independent of $n$ and $L$.
\end{enumerate}
\end{theorem}

\begin{proof}
(i) Let $\mathcal K=L^{-4}\sum_x\mathcal P_x$ and let $h$ be gauge invariant at
every configuration with $\Lip(h)<\infty$, and let $U,U'$ differ only at $i$.
If $i\in V_x$, then $\mathcal P_xh(U)=\mathcal P_xh(U')$. If $i\notin V_x$,
Lemma~\ref{lem:softlip} and any coupling $\pi$ of $\gamma^0_{V_x,U}$ and
$\gamma^0_{V_x,U'}$ give $|\mathcal P_xh(U)-\mathcal P_xh(U')|\le\Lip(h)(\int
D\,d\pi+\delta(U_i,U'_i))$, and the infimum over $\pi$ gives
$\Lip(h)(\tilde R_x(i)+1)\delta(U_i,U'_i)$. As $i$ lies in $n^4$ cubes,
Proposition~\ref{prop:starresponse} gives $\Lip(\mathcal Kh)\le(1-\kappa)
\Lip(h)$ with $\kappa=(1-259/n)n^4/L^4>0$, and $\mathcal Kh$ is again gauge
invariant at every configuration, by gauge covariance of $\gamma^0_{V_x,U}$.
The spectral argument of Proposition~\ref{prop:starcond} gives
$\ip h{(I-\mathcal K)h}\ge\kappa\Var_\mu(h)$, first for such $h$ and then, by
density, for every real $f\in\mathcal A$; and $L^{-4}\sum_x\mu((f-\mathcal
P_xf)^2)=\ip f{(I-\mathcal K)f}$, which is $\BP(n,1-259/n)$.

(ii) For $n\ge518$ use (i). For $2\le n\le517$ use the proof of
Theorem~\ref{thm:stars}, which gives $\BP(n,(1-6c\,s(c))(n-1)^4/(n^4\lceil
(n-1)^4/2\rceil))$ whenever $6c\,s(c)<1$; here $1-6c\,s(c)\ge1-\frac{1827}
{1933}=\frac{106}{1933}$ by Lemma~\ref{lem:svalues}(ii), and $(n-1)^4\ge
\lceil(n-1)^4/2\rceil$, so $\BP(n,\frac{106}{1933}n^{-4})$ holds, and
$\frac{106}{1933}n^{-4}\ge\varpi$.
\end{proof}

Theorem~\ref{thm:starsuniform} extends the uniform range $\frac5{18}$ of
Theorem~\ref{thm:bpweighted} past $\beta_W=\frac13$, where
Proposition~\ref{prop:methodlimit} stops every gauge-fixed single-link
scheme and Proposition~\ref{prop:facebarrier} every scheme that treats face
links as single sites. The faces cost nothing here; the margin is lost at the
cube edges, through the column $h(0,3)$ of Lemma~\ref{lem:column}.\footnote{A
linear programme over weights depending on the block, with the exact
constants for $z(v)=2$, suggests on cubes of side up to $8$ that the method
reaches the star threshold $\beta_W\approx0.3609$ of Theorem~\ref{thm:stars};
with weights depending only on $z(v)$ it stops near $\beta_W\approx0.345$.
These are numerical observations and no statement rests on them.} The value
$\beta_W\le\frac7{20}$ is far below the scaling region of four-dimensional
$\SU(2)$, where $\beta_W$ is of order $2$ and above, and the next proposition
explains why no one-step contraction of this kind reaches it.

\subsection{Single links near flat configurations}

Let $A(\kappa)=\int x_1\,\nu_{\kappa e_1}(dx)$ be the mean resultant length of
$\nu_b$ at $|b|=\kappa$. By symmetry $\int x\,\nu_b(dx)=A(|b|)b/|b|$; $A$ is
nondecreasing, its derivative being the variance of $x_1$, and
$A(\kappa)\to1$ as $\kappa\to\infty$.

\begin{proposition}\label{prop:flat}
Let $\beta>0$, let $e$ be a link whose six plaquettes contain no link fixed at a
value other than $1$, and let $k$ be a partner of $e$. The single-link influence
satisfies
\[
C_{ek}=\sup\frac{W_1(\nu_b,\nu_{b'})}{\delta(U_k,U'_k)}\ge\frac{A(24\beta)}6 .
\]
Hence in any single-link scheme in which every link has at least twelve free
partners, as for the forest of Theorem~\ref{thm:A}, every row sum of the
influence matrix is at least $2A(24\beta)$, which tends to $2$ as
$\beta\to\infty$.
\end{proposition}

\begin{proof}
Take all links equal to $1$, so every staple is $1$ and $b=24\beta\cdot1$.
Change $U_k$ so that the staple of the plaquette containing $e$ and $k$
becomes $\cos\vartheta+t\sin\vartheta$ with a unit imaginary quaternion $t$;
then $\delta(U_k,U'_k)=2\sin\frac\vartheta2$ and $b'=4\beta(5+\cos\vartheta+
t\sin\vartheta)$. The function $x\mapsto\ip xt$ is $1$-Lipschitz, so
$W_1(\nu_b,\nu_{b'})\ge\ip{\int x\,d\nu_{b'}-\int x\,d\nu_b}t=A(|b'|)\,
4\beta\sin\vartheta/|b'|$. Divide by $2\sin\frac\vartheta2$ and let
$\vartheta\to0$. For the limit of $A$: the law of $x_1$ under
$\nu_{\kappa e_1}$ has density proportional to $e^{\kappa t}\sqrt{1-t^2}$ on
$[-1,1]$, and for $0<\epsilon<\frac12$ the mass of $[-1,1-2\epsilon]$ divided
by that of $[1-\epsilon,1]$ is at most $e^{-\kappa\epsilon}\pi/\int_{1-
\epsilon}^1\sqrt{1-t^2}\,dt\to0$.
\end{proof}

At weak coupling the Gibbs measure concentrates near flat configurations, so
even an influence matrix computed only on typical configurations does not
contract there: per link it approaches $\frac16$ per partner, and
$12\cdot\frac16=2$.\footnote{Numerically $A(\kappa)=\frac12$ at
$\kappa\approx2.45$, that is $\beta_W\approx0.41$.} Larger blocks with better
internal estimates push the threshold to some finite $\beta_W$, but the mass at
weak coupling is non-perturbative, $m\sim\Lambda e^{-c\beta}$, and a
contraction can appear only at block sizes of order $e^{c\beta}$ lattice units.
Section~\ref{sec:weak} and Appendix~\ref{app:obstructions} make this precise.
